\section{Artifact layout, public APIs, and reproduction}
\label{main:reproduction}

The accompanying archive contains the complete article sources and PDF,
an additive integration patch, a runnable source snapshot, fixed source
corpora, proof examples, experiment records, and validation logs.  The snapshot
is included so that local reproduction does not depend on the current state
of the repository's main branch.  The patch is separately based on the exact
revision printed on the title page.

\subsection{Delivered files and their roles}

\begin{center}\small
\begin{tabularx}{\textwidth}{@{}p{.32\textwidth}Y@{}}
\toprule
Path & Purpose\\
\midrule
\path{article.tex}, \path{sections/} & Complete editable article sources;
the bibliography is inside the main source.\\
\path{article.pdf}, \path{build.sh} & Typeset article and local rebuilding command.\\
\path{integration.patch} & New modules, tests, and reproduction data for the
pinned repository tree; no existing runtime file is overwritten.\\
\path{runtime/fast/} & Standard-library native implementation, its pinned
dependencies, four new test modules, and fixture/benchmark tools.\\
\path{runtime/synthesis/}, \path{runtime/reports/57/} & Source corpora needed
by the literal boundary and retained-sector audits.\\
\path{examples/} & Small independently verified source-bound proof packets,
including the Q/standard mixed-ray example.\\
\path{results/} & Case-level audit and timing JSON, source hashes, and logs.\\
\path{tools/reproduce.py} & Portable smoke, focused-test, audit, and benchmark
entry points.\\
\path{provenance/} & Pinned revision, file hashes, environment, patch replay
checks, and PDF verification record.\\
\bottomrule
\end{tabularx}
\end{center}

The native runtime requires Python and its standard library.  The recorded
independent explicit-cut comparison uses Regina 7.4.  A machine without Regina
can run the new native tests and replay the supplied small certificates;
the explicit comparison mode requires Regina.  Rebuilding the paper requires
a LaTeX installation with the packages named in \path{article.tex}.  The
included upstream runtime license is preserved.

\subsection{A source-bound whole-cut example}

From \path{runtime/fast/}, the following computes and independently checks
the cut along a supplied primitive meridian:

\begin{lstlisting}[language=Python]
from fastunknot.normal_cut_handlebodies import normal_cut_handlebodies
from fastunknot.normal_cut_handlebodies_verify import (
    verify_normal_cut_handlebodies,
)
from normal_orbit_research.fixtures import layered_torus

source, coordinates = layered_torus(4)
answer = normal_cut_handlebodies(
    source, coordinates, record_certificate=True,
)
assert answer['status'] == 'COMPLETE'
assert verify_normal_cut_handlebodies(
    source, coordinates, answer['certificate'],
)
assert answer['census']['histogram'] == [
    {'genus': 0, 'multiplicity': 1},
]
skeleton = answer['certificate']['boundary']['skeleton']
assert len(skeleton['ports']) == 15
assert skeleton['interval_attachments'] == 3
\end{lstlisting}

This is a theorem about the supplied cut.  Converting it to a knot verdict
requires the additional surface and exterior hypotheses described above.
An unresolved result should be passed to a different exact method or retained
as unresolved; it must not be converted to a negative recognition answer.

For a large common multiplicity, multiply the normal coordinates by the exact
integer.  The native algorithms consume those binary integers directly.
When saving a large certificate in JSON, use the established exact codec:

\begin{lstlisting}[language=Python]
import json
from fastunknot.integer_codec import json_safe

source, primitive = layered_torus(1)
m = 1 << 16384
coordinates = [[m * n for n in row] for row in primitive]
answer = normal_cut_handlebodies(
    source, coordinates, record_certificate=True,
)
assert answer['census']['ball_components'] == m
wire = json.dumps(json_safe(answer['certificate']))
proof = json.loads(wire)
assert verify_normal_cut_handlebodies(
    source, json_safe(coordinates), proof,
)
\end{lstlisting}

The codec uses explicitly signed hexadecimal strings for sufficiently large
integer fields, avoiding process-wide decimal-conversion changes.  Equal
Boolean and integer values are not interchangeable schema fields.  The checker
validates the prescribed types and source correspondence; successful JSON
parsing is not itself certificate validation.

\subsection{A sector example that must retain the mixed standard ray}

\begin{lstlisting}[language=Python]
from sector_block_research.fixtures import capped_bipyramid_chain
from fastunknot.sector_blocks import certify_sector_block_rays
from fastunknot.sector_blocks_verify import verify_block_ray_certificate

fixture = capped_bipyramid_chain(2)
source = fixture['triangulation']
support = fixture['allowed_types']
q = certify_sector_block_rays(source, support, phase='quadrilateral')
s = certify_sector_block_rays(source, support, phase='standard')
assert len(q['rays']) == 4
assert len(s['rays']) == 5
assert verify_block_ray_certificate(source, q['certificate'])
assert verify_block_ray_certificate(source, s['certificate'])
\end{lstlisting}

The source and support are the same.  The difference in ray counts is required
by the proved canonical-minimum interaction, rather than a tolerance or a
heuristic enumeration choice.  The negative canonical-Euler and positive
canonical-Euler certificates have separate schema names and weaker topological
meaning than an essential-disk certificate.

\subsection{Commands and scope of the validation records}

From the archive root:
\begin{lstlisting}
python tools/reproduce.py smoke
python tools/reproduce.py tests
python tools/reproduce.py audits
python tools/reproduce.py benchmarks
bash build.sh
\end{lstlisting}

The smoke command verifies the saved small proof packets.  The test command
runs the four new test modules; it does not claim to recreate all 1,644 tests
of the integrated source tree.  The audits use the included fixed source
corpora and write new reports under \path{reproduced/}.  Benchmark commands
run the recorded algorithms with their original sample counts; wall-clock
times will depend on the host.  The binary capacity cases report exact
arithmetic and certificate volume, not an explicit-oracle speed ratio.

To reproduce the full integrated suite, check out the pinned repository,
inspect and apply \path{integration.patch} at the repository root, enter
\path{Topology/UnknotRecognition/fast/}, and run:
\begin{lstlisting}
python -m unittest discover -s tests -v
\end{lstlisting}
The delivered log records the actual 1,644-test run.  A subsequent focused
sector run validates the final generalized fixture as well.  The separate
runtime snapshot is sufficient for this contribution's focused tests and
audits; it intentionally does not duplicate every unrelated report in the
large ProveIt repository.

The integration patch only creates new files.  It does not enable a new
diagram recognizer policy, replace existing source checkers, or assume that
all future normal cuts will satisfy the disk-interface criterion.  Its
purpose is to make the proven local operations and their limits reviewable
before a hierarchy consumer depends on them.

\subsection{Results that are proved but not implemented in this delivery}

Four extensions are deliberately separated from the measured code:
the three-unweighted-count partition certificate, reuse of a primitive
certificate under the replication theorem, exhaustive bounded-genus survivor
search, and disk-base full-annulus absorption.  Their mathematical statements
and proofs are in the article.  The current production APIs do not silently
claim their speed, coverage, or proof format.  The annular eligibility
diagnostic tests only the stated base-type preconditions.

This distinction makes a later implementation comparison precise.  A future
patch can add one extension, replay the same source bank, and identify exactly
which proof language, accepted cut class, and measured work changed.
